\section{Introduction}\label{sec:introduction}

\subsection{Motivation}
\todo{Driving as a sequence of decisions under uncertainty (lane changes, overtaking, speed adaptation);
the trade-off between safety and efficiency; why RL is a natural framework.}

\subsection{State of the art}\label{sec:sota-intro}
\todo{Draft and extend.}
\begin{itemize}
  \item \textbf{DRL for autonomous driving:} surveyed in~\cite{kiran2022drl}. \todo{perception vs.\ decision-making, typical tasks.}
  \item \textbf{HighwayEnv as a benchmark:} \env{HighwayEnv}~\cite{highway-env} is a standard lightweight
        environment for tactical decision-making on highways.
  \item \textbf{Simulation-to-reality gap:} policies trained in simplified simulators may not transfer to real
        traffic~\cite{zhao2020sim2real}. \todo{what this implies for the conclusions of this study.}
\end{itemize}

\subsection{Contributions}
\todo{Bullet list of what was actually done: two from-scratch implementations, shared evaluation protocol,
reward study, density study, \ldots}

\subsection{Report structure}
\Cref{sec:problem} formulates the MDP, \cref{sec:methods} presents the baseline and the two algorithms,
\cref{sec:reward} the reward designs, and \cref{sec:protocol} the experimental protocol.
\Cref{sec:results} reports the results, \cref{sec:discussion} discusses them, and
\cref{sec:conclusion} concludes. The appendix gives reproduction instructions and full hyper-parameter tables.
